\scanpage{40}
% L1–3
semi-exact :
\[
\left\{\begin{aligned}
 &\text{Si }T^i(B)=0,\ \text{alors }T^i(A)=0\ \text{pour tt facteur direct }A\text{ de }B\\
 &\text{Si } A\text{ a une filtration finie à facteurs }A_\alpha\\
 &\qquad\text{tels que }T^{\unsure{i|v}}(A_\alpha)=0\ \text{pour tt }\alpha,\ T^i(A)=0
\end{aligned}\right.
\]

% L4–12
Pour les foncteurs cohomologiques de faisceaux de $l$-torsion \uline{constructibles} sur $X$ (\unsure{pas qct pas séparé}) ces remarques \deleted{déjà} impliquent déjà que pour avoir $T^i(F)=0$ pour tt tel $F$, il suffit de le prouver pour un $F$ de la forme $i_!(p_*(\unsure{\mathbb Z|Z}/l\mathbb Z)_{Y'})$\inserted{$=q_*j_!(\mathbb Z/l\mathbb Z)_{Y'})$}, où $Y$ est un sous-préschéma constructible de $X$, $Y'$ un rev. étale fini de $Y$ de degré constant, et $p:Y'\unsure{\mathord{-}|\to}Y$ et $i:Y\to X$ les morphismes can. On \unsure{pt|vt} prendre $j$ une immers. ouverte dans $X'$ fini sur $X$.

% L12–14, left margin
\diagram{\begin{tikzcd}[ampersand replacement=\&,row sep=small,column sep=small]
Y'\arrow[r,"j"]\arrow[d,"p"']\&X'\arrow[d,"q"]\\
Y\arrow[r,"i"']\&X
\end{tikzcd}}

% L13–21
Toute sous-catégorie pleine de la catégorie des faisceaux de $l$-torsion constructibles sur $X$, stable par facteurs directs, par extension, et contenant les faisceaux $q_*(j_!(\mathbb Z/l\mathbb Z)_{Y'})$, précédents, est toute // Si par ailleurs cette catégorie est \deleted{\unsure{aussi}} telle que pour tte suite exacte $0-A-B-C-0$, si $B$ et $C$ y sont, alors $A$ y est --- p. ex si c'est une sous-catégorie \uline{épaisse} --- alors il suffit de vérifier que les $q_*(\mathbb Z/l\mathbb Z)_{X'}$ y sont, $X'$ fini sur $X$.

% L22–24
Si $X$ \unsure{noeth}, on peut même \unsure{tj se|tjrs} prendre $X'$ intègre. Si $X$ \deleted{\unsure{est de caractéris} \illegible{1}} \inserted{de plus} univ. japonais, on peut prendre $X'$ intègre normalisé.

% archive sheet number: 39
